\section{Why the proof belongs to iPRA}
\label{sec:ipra}

The mathematical argument above uses only rational arithmetic and finite
combinatorics. That observation alone would not establish a weak-theory
result: a true statement about a terminating computation can require a
stronger theory to prove uniformly. We now specify the definitions and
induction invariants that justify the iPRA claim.

\subsection{The underlying theory}

Use ordinary first-order intuitionistic logic with equality and one
function symbol for each primitive recursive definition. The arithmetic
axioms are the defining equations for zero, successor, projections,
composition and primitive recursion, successor injectivity,
\(Sx\ne0\), and induction restricted to quantifier-free formulas.
Parameters are allowed in the induction schema.

This is exactly the presentation defined by \code{PRAAx} in the inspected
ConciseIPRA source \cite{proveit}; \code{PRA.qfInd} exposes its induction
rule. Thus we do not silently replace the requested theory by Heyting
arithmetic with induction for all formulas.

\subsection{Arithmetic and rational codes}

Addition and multiplication are defined by primitive recursion.
Truncated subtraction, comparison, parity, bounded quotient and remainder,
finite sums, products, factorial, and exponentiation are primitive recursive.
The usual ordered semiring identities follow by quantifier-free induction
on their defining recursions. For instance, distributivity is proved by
induction on the recursion argument of multiplication, with the other
arguments as free parameters.

An integer can be represented by a pair \((a,b)\) denoting \(a-b\).
Equality is \(a+d=c+b\) for \((a,b)\) and \((c,d)\).
Addition and multiplication have polynomial definitions:
\[
 (a,b)+(c,d)=(a+c,b+d),\qquad
 (a,b)(c,d)=(ac+bd,ad+bc).
\]
Negation swaps the entries. Order is expressed by a natural-number
comparison. This permits all signed calculations involving \(\mu\),
\(\lambda_d\), and the errors \(E_v\) without adjoining negative numbers
as a new domain.

A rational is an integer numerator code and a positive natural denominator.
Reduction to lowest terms is unnecessary. For denominators \(b,d>0\),
\[
 \frac ab+\frac cd=\frac{ad+bc}{bd},\qquad
 \frac ab\frac cd=\frac{ac}{bd},\qquad
 \frac ab\le\frac cd\ \Longleftrightarrow\ ad\le cb.
\]
Inverting a nonzero rational swaps the absolute numerator and denominator
and retains its sign. The proof uses inverses only with a separately proved
positive input. In particular \(g(d)>0\), \(h(d)>0\), \(G(z)\ge1\);
the factors \(p-\rho(p)\) are positive.
All denominator clearing is therefore legitimate in the integer coding.
Rational squares are nonnegative by a decidable sign comparison followed
by the integer multiplication identities.

Finite rational sums are recursive folds of these explicit operations.
Their algebraic rules are proved by induction on fold length.
For example, double-sum interchange is obtained by induction on one
length and then on the other, using distributivity and associativity at
each step. Nonnegative-term comparison is a fold induction on the
partial-sum inequality. These induction formulas compare particular
integer codes, and are quantifier-free.

\subsection{Bounded conditions can be put into induction invariants}

For a primitive recursive Boolean predicate \(P(i,\vec a)\), define
\[
 \operatorname{All}_P(0,\vec a)=1,\qquad
 \operatorname{All}_P(n+1,\vec a)
   =\operatorname{All}_P(n,\vec a)\cdot P(n,\vec a).
\]
Its value is \(1\) exactly when all entries \(i<n\) pass the test.
The analogous bounded existential test is obtained by testing whether the
sum of the Boolean entries is positive. Bounded searches and bounded
maxima are defined by folds with a success flag. For an existential test,
the fold also records the first successful index. Its invariant says that
if the flag is \(1\), the stored index is below the processed bound and
passes \(P\). This supplies the witness constructively.

The facts required to use these encodings are themselves proved by
quantifier-free induction. A useful explicit strengthened invariant for
recovering entries is
\[
 \operatorname{All}_P(n,\vec a)=1\ \wedge\ i<n
       \ \Longrightarrow\ P(i,\vec a)=1,
\]
with \(i,\vec a\) as free parameters; the induction step splits the
decidable cases \(i<n\), \(i=n\), and \(i>n\).
In the other direction, under assumptions supplying the bounded entries,
the defining multiplication recurrence proves that the aggregate is \(1\):
with the final bound \(b\) fixed, induct on the quantifier-free assertion
\(j\le b\Rightarrow\operatorname{All}_P(j,\vec a)=1\).
This establishes the familiar bounded-quantifier notation without
induction on an unbounded universal formula.

For a nested bounded formula, apply this construction from the innermost
test outward. Its characteristic function remains primitive recursive,
and its use inside an induction can be replaced by equality of that
function to \(1\). The finite assertions in the present proof involve
only such bounded conditions and coded arithmetic equalities or
inequalities.

\subsection{Explicit bounds for the finite objects}

Here are bounds sufficient to make the construction literal.
They need not be efficient.
\begin{itemize}
\item The least-divisor search on \(n>1\) tests at most \(n\) candidates.
Factorization at \(0\) is assigned a fixed default, and its correctness
claims concern positive inputs. The factorization loop has \(n\) stages.
A padded factor trace has length
\(n\), and entries at most \(n\); it can be stored as digits in base \(n+1\),
together with the length \(\ell\le n\) of its active prefix. Explicitly,
the code is \(\operatorname{digits}+\ell(n+1)^n\), with
\(\operatorname{digits}<(n+1)^n\), hence is below \((n+1)^{n+1}\).
\item The subtraction Euclidean algorithm has at most \(a+b\) active
stages. Its Bezout coefficients start with absolute value at most \(1\)
and at each stage are copied or subtracted. Their absolute values are
bounded by \(2^{a+b+1}\). Sign and magnitude codes give a bounded trace.
\item Residue tables modulo \(v\) have \(v\) entries. In the sieve,
\(d,e\le z\) and \(\lcm(d,e)\le z^2\), so all such tables have at most
\(z^2\) entries. Their truth values are computed by remainder tests.
\item Divisor and squarefree tests at \(d\le z\) use ranges bounded by
\(z\). The sums defining \(G,W_d,\lambda_d\) have at most \(z\) terms.
The quadratic form has at most \(z^2\) terms.
\item In the harmonic comparison, \(ab\le s^2\).
For each \(d\le s^2\), there are at most \(s^2\) prime slots and
two exponents in \(\{0,\ldots,s\}\) at each slot.
Thus \((s+1)^{2s^2}\) bounds the number of padded exponent codes.
Their products and weights are computed by bounded folds.
\item A dyadic shell index for the tail below \(L\) is found in the
explicit range \(m,\ldots,m+L\). No logarithm function or unbounded
least-witness principle is required.
\end{itemize}
The coding of an entire rational computation is likewise a primitive
recursive trace: each next numerator and denominator is computed by the
displayed operations. If a bound on trace entries is wanted in an
invariant, compute the maximum along that same finite trace by a fold.
There is no existentially postulated sequence of intermediate rationals.

\subsection{The induction audit}

Every induction used in the proof can be organized as one of the following
finite invariant checks. Ambient bounds such as \(n,z,s,X,m,L\) are fixed
parameters during the relevant induction.

\renewcommand{\arraystretch}{1.12}
\begin{longtable}{@{}p{0.21\textwidth}p{0.28\textwidth}p{0.44\textwidth}@{}}
\toprule
Part of proof & Induction variable & Quantifier-free invariant after coding\\
\midrule
\endhead
Division & Number of processed increments &
Dividend \(=\) divisor \(\times\) quotient \(+\) remainder, with remainder
below divisor.\\
Factorization & Bounded stage counter &
Original number \(=\) remaining factor \(\times\) recorded product;
the recorded factors pass a bounded primality test; active steps decrease
the remaining factor.\\
Bezout and gcd & Bounded stage counter &
The two stored coordinates have the stored linear-combination
representations; common divisors up to \(a+b\) agree with the initial
pair; either the computation has stopped or its sum has decreased.\\
Factor uniqueness & Length of cancelled prefix &
Matching prime factors have been removed, and the two remaining
products agree; all residual factors are prime.\\
Chinese remainders & Length of prime prefix &
The computed reconstruction and residue projection are inverse on their
bounded domains; the product of root counts equals the root-table count.\\
Mobius cancellation & Length of prime prefix &
The partial subset sum equals the corresponding product of local
\((1-1)\) factors.\\
Finite algebra & Sum/product length &
The two accumulated rational codes are equal, or one accumulated code
is at most the other.\\
Selberg transform & Divisor-sum length &
The accumulated reindexed sums agree; the interval cancellation identity
holds for indices bounded by \(z\).\\
Quadratic identity & Bounded sum lengths &
The partial quadratic expansion equals its reindexed divisor expression;
the completed expansion equals \(1/G\).\\
Harmonic comparison & Prime-prefix and box counters &
Partial exponent-box sums factor; original pair codes map injectively to
box codes; all weights are nonnegative.\\
Dyadic estimates & \(t\) or shell counter &
\(H_{2^t}\ge1+t/2\), \((t+1)^2\le4^t\), or the displayed finite
telescoping inequality.\\
\bottomrule
\end{longtable}

For the termination statements in this table, a numerical strengthening
removes any unbounded termination assertion. For the factor loop at stage
\(i\le n\), either it has reached \(1\), or its remaining factor plus \(i\)
is at most \(n\). At stage \(n\) the latter case contradicts positivity.
For the subtraction Euclidean loop, either it has stopped, or the
coordinate sum plus the stage counter is at most the initial \(a+b\).
At stage \(a+b\) two positive coordinates are impossible.
Both strengthened assertions are quantifier-free after the state is coded.

In the factor-uniqueness and CRT rows, statements about all entries are
bounded by the fixed input products or trace lengths. Their bounded
Boolean aggregates can be carried in the invariant. The elementary
preservation calculations are those already given in
Section~\ref{sec:finite}; no induction over all arbitrary formulas or all
unbounded witnesses is needed.

\subsection{The resulting derivability claim}

Expand the finite definitions using their primitive recursive symbols,
interpret integer and rational notation by the preceding codes, and
replace bounded conditions in induction invariants by their characteristic
functions. The proofs in Sections~\ref{sec:finite}--\ref{sec:convergence}
then use the defining equations, intuitionistic reasoning with decidable
arithmetic cases, and the quantifier-free inductions just listed.
After the final denominator multiplication they derive
\eqref{eq:ipra-target}; universal introduction closes the free parameters
\(q,U,V\).

Only the finitely many specific primitive recursive definitions used by
this proof are involved. No primitive recursive universal interpreter for
all primitive recursive programs is assumed.

This is a human-readable proof of the iPRA theorem, together with its
translation discipline. It is not a claim that a formal derivation
certificate has already been generated. Producing that certificate is
the object-theory stage of the Lean plan below.

In particular, the proof requires neither an infinite Euler product,
a prime-distribution asymptotic, a compactness or completeness principle,
nor a choice of an unknown convergence threshold. All numerical case
splits concern decidable tests. There is no use of excluded middle for an
arbitrary proposition, Markov's principle, or a decision about infinitely
many twin primes.
